\section{결론}
\label{sec:conclusion}

추론 보강형 E2E 주행 데이터는 행동 이유를 제공하지만, 반복 CoC 사건이 언제 후보 의무를 열고 자차 움직임 응답이 언제 이를 닫는지, 그리고 그 연결이 어떤 출처에 근거하는지는 별도 감사가 필요하다. 본 논문은 시점이 있는 CoC 주석, 도출 자차 움직임 응답, 반복ㆍ기한 정책과 장면 출처를 \uppaal 시간 관찰자 네트워크로 컴파일하여 저장된 사건열의 시간ㆍ출처 계약을 검사하는 방법을 제안하였다.

94개 장면의 단일 이벤트 선별에서 얻은 134개 방향 명시 주석 사건별 후보 계약은 $D=3$초 분석 설정에서 125개 속도 변화 계약 통과, 4개 이미 만족된 계약 통과, 5개 검토 후보로 분기되었다. 반복 사건은 응답 시작을 공유할 수 있고 상태형 분석은 선택 사례에만 적용했으며, 0.5~m/s 이미 충족 임계값도 코호트 보정 탐색 논리이다. 따라서 결과를 안전/위험 레이블이 아니라 약한 유지 후보, 정책ㆍ기한ㆍ탐지기ㆍ계보별 검토 후보 및 출처 불충분 항목으로 유형화했다. 별도 파서ㆍ탐지기의 13체인 수치는 구현 파이프라인 스모크 출력으로만 보존하며 공유 계약의 정량 결론에서 제외한다. 집계 산출물의 13/13 기록도 실행 가능한 기준선과 항목별 출력이 없어 독립 재현성 근거로 사용하지 않는다.

이 결과는 물리적 안전성, CoC의 일반적 충실도 또는 제안 라우팅에 따른 학습 성능 향상을 입증하지 않는다. 더 강한 결론을 위해서는 공식 출처가 확인된 장기 시퀀스 데이터, CoC와 독립적인 궤적ㆍ응답 검증기, 그리고 학습 정책의 폐루프 평가를 후속 증거 계층으로 결합해야 한다.
